\chapter{Conclusions}

This thesis established the fundamental structure of quantum error-correcting codes. By focusing on mathematical structure and algebraic consistency, it addressed the limitations imposed by quantum theory, such as the no-cloning theorem and the continuous nature of errors. In doing so, it demonstrated the efficacy of encoding information into larger Hilbert subspaces to protect it from environmental decoherence.

First, defining and treating errors as quantum operations over the system and its environment laid the necessary groundwork. In particular, it was explained why the focus is placed on independent qubit errors and how any continuous error can be decomposed into a basis of Pauli operations, ensuring that correcting these discrete errors automatically corrects the continuous ones. Building upon this foundation, the core structure of a quantum code was defined, introducing the Knill-Laflamme conditions as the definitive mathematical criteria for successful error recovery. Additionally, key theoretical concepts, such as the quantum Hamming bound and degenerate codes, were explored.

The core of this work centered on the stabilizer formalism and its corresponding codes, which serve as the primary framework for efficiently constructing and analyzing quantum error-correcting codes. To this end, it was initially necessary to examine classical linear codes, which act as the classical analogue to stabilizer codes, particularly regarding the use of parity-check matrices for efficient error detection. Following this, the study of the symplectic group and the derivation of extended symplectic operations—which can be systematically constructed in quantum circuits—demonstrated how this algebraic structure functions analogously to classical parity checks, enabling efficient syndrome measurement.

Finally, the theoretical framework was applied to two concrete examples. First, the quantum repetition code was analyzed using the Knill-Laflamme conditions, evaluating its fidelity and illustrating its inherent inability to correct phase-flip errors. Second, the general Shor code was examined, successfully deriving its necessary stabilizers to prove that it functions as a valid stabilizer code capable of correcting arbitrary single-qubit errors.

In conclusion, the rigorous development of these theoretical foundations provides a solid understanding of quantum error-correcting codes. Moreover, it equips the reader with the necessary analytical tools to design and construct new fault-tolerant codes utilizing the stabilizer formalism.
